% Required packages (add to main.tex preamble if not already present):
%   \usepackage{booktabs}

\section{Experimental Setup}
\label{sec:setup}

\subsection{Datasets}

We adopt four classification benchmarks from GLUE~\cite{wang2018glue}, chosen to cover diverse
linguistic phenomena and to stress-test forgetting across syntactically and semantically distinct
tasks.

\paragraph{SST-2} (Stanford Sentiment Treebank~\cite{socher2013sst}) is a binary sentiment
classification task over single movie-review sentences.
The full training set contains 67{,}349 examples; we report accuracy on the official validation
split (872 examples).

\paragraph{MRPC} (Microsoft Research Paraphrase Corpus~\cite{dolan2005mrpc}) is a paraphrase
detection task over sentence pairs drawn from online news articles.
The full training set contains 3{,}668 sentence pairs; we report macro-F1 on the validation
split (408 examples).

\paragraph{CoLA} (Corpus of Linguistic Acceptability~\cite{warstadt2019cola}) requires binary
classification of whether an English sentence is grammatically acceptable.
The full training set contains 8{,}551 examples; we report Matthews Correlation Coefficient (MCC)
on the validation split (1{,}043 examples).

\paragraph{MNLI} (Multi-Genre Natural Language Inference~\cite{williams2018mnli}) is a three-class
entailment task (entailment / neutral / contradiction) over premise--hypothesis pairs spanning
multiple genres.
The full training set contains 392{,}702 examples; we use the \emph{matched} validation split for
evaluation and report accuracy.

For all tasks we uniformly cap training at $N_{\mathrm{train}} = 2{,}000$ examples per task to
ensure each full continual learning run completes within a single Colab T4 session; the cap is
binding only for SST-2 and MNLI, while MRPC and CoLA use their full training sets.
The task sequence order is fixed as SST-2\,$\to$\,MRPC\,$\to$\,CoLA\,$\to$\,MNLI ($K = 4$).

\subsection{Models}

We evaluate two instruction-tuned small LLMs.

\paragraph{Qwen2.5-1.5B-Instruct}~\cite{qwen2024} is a 1.5-billion parameter chat model from the
Qwen2.5 family.
Its small footprint permits multiple experimental runs---including full ablation sweeps---within the
Colab memory budget, making it the primary model for all six methods and all ablations.

\paragraph{Phi-3.5-mini-instruct}~\cite{abdin2024phi3} is a 3.8-billion parameter model from
Microsoft's Phi series.
We include it to test whether findings generalize beyond the Qwen architecture, running only the
main six-method comparison (not ablations) due to its larger memory footprint.

Both models use instruction-following chat templates, enabling uniform prompt-based task formatting
across methods.
Both are released under permissive licenses (Apache~2.0 and MIT, respectively) that explicitly
permit fine-tuning for research purposes.

\subsection{Hyperparameters}

Table~\ref{tab:hparams} lists all fixed hyperparameters used across every method and seed.
Method-specific coefficients ($\lambda$, $B$, $\mu$, $f$) were set to literature-consistent
defaults; we examine the sensitivity of MMR to its replay ratio $f$ in
Section~\ref{sec:results}.

\begin{table}[h]
\centering
\caption{Hyperparameters used across all experiments.}
\label{tab:hparams}
\begin{tabular}{ll}
\toprule
\textbf{Hyperparameter}               & \textbf{Value}                \\
\midrule
Learning rate                         & $2 \times 10^{-4}$            \\
Effective batch size                  & 8                             \\
Maximum sequence length               & 256 tokens                    \\
Epochs per task                       & 2                             \\
Optimizer                             & \texttt{paged\_adamw\_8bit}   \\
LoRA rank $r$                         & 16                            \\
LoRA scaling $\alpha$                 & 32                            \\
LoRA dropout                          & 0.05                          \\
Target modules                        & \texttt{q,k,v,o\_proj}        \\
EWC coefficient $\lambda$             & $500$                         \\
Replay buffer size $B$ (per task)     & $200$ examples                \\
O-LoRA orthogonality coeff.\ $\mu$    & $0.1$                         \\
MMR replay ratio $f$                  & $0.06$                        \\
Random seeds                          & 42, 123, 2025                 \\
\bottomrule
\end{tabular}
\end{table}

\subsection{Hardware and Reproducibility}

All experiments are run on an NVIDIA T4 GPU (16\,GB VRAM) via Google Colab free tier.
We fix random seeds $\{42, 123, 2025\}$ and report means and standard deviations across the three
runs.
Wall-clock time is measured with \texttt{time.time()} surrounding the training loop; peak GPU
memory is recorded with \texttt{torch.cuda.max\_memory\_allocated()} immediately after training.
Code, prompts, and result JSON files are available at \texttt{https://github.com/gautamdubey-ethara/RPaper} (access on request).
